\subsection{An analytically solved policy benchmark}\label{sec:balanced-policy}

The full policy game can also be solved without choosing
a numerical policy grid under an economically interpretable
balance of technological channels: $\rho=\eta$.
Technical compatibility then raises willingness to pay
through network value by exactly the amount needed
to offset its effect on the \emph{fulfilled-expectations}
Cournot matrix. Firms still optimize residual differentiation.

\begin{proposition}[Maximum-depth international standards when
the two technological channels balance]\label{prop:analytic-D}
Let $\rho=\eta=r\in(0,1)$ and $\beta=1-r$,
$\alpha=2-r$, $d>0$, and suppose
\begin{equation}\label{eq:analytic-D-conditions}
 \kappa>
 \max\left\{
 \frac{54\beta^2d^2}{r^2\alpha^2},
 \frac{6\beta d^2}{\alpha^3}
 \right\}.
\end{equation}
Then the privately chosen design Nash is unique,
symmetric, and interior after every governmental history,
and neither designs nor quantities depend on the
government's standardization regime or depth.
With $z$ denoting the optimized residual homogeneity
and $Q=d/(\alpha+2\beta z)$, country payoffs satisfy
\begin{equation}\label{eq:analytic-D-welfare}
 m(g)=o(g)=s+r gQ^2,\qquad
 i(g)=s+3r gQ^2.
\end{equation}
The specified bilateral-first government game selects
the maximum compatibility depth $g^*=1$ and a substantive
international standards union $\IS$ on the equilibrium path.
\end{proposition}

\begin{proof}
If $\delta=\rho-\eta=0$, the Cournot matrix
$M=I+\beta X$ is independent of the governmental
technical-link matrix $H$ and depth $g$.
The first bound in \eqref{eq:analytic-D-conditions}
guarantees full-asymmetric strict diagonal concavity
of the firms' design game and hence a unique Nash.
The second ensures the symmetric unilateral FOC
is positive at zero; it is strictly negative at one,
yielding an interior symmetric stationary point,
which must be that unique Nash.
The single AB link adds $rgQ^2$ to domestic
consumer surplus; an international union contains
three such links. Private profit and design costs
are identical in every regime. Equation
\eqref{eq:analytic-D-welfare} follows, yielding
$i(g)>m(g)=o(g)>s$ for every $g>0$.
The outsider applies, both incumbents approve,
and the initially invited member accepts.
Since $i(g)$ is strictly increasing, proposer
$A$ optimally chooses $g^*=1$.
\end{proof}

This benchmark identifies a condition under which
the government obtains full standardization while
the design margin remains \emph{endogenous}.
It is a balanced-channel special case, not a
general threshold separating international
coordination from fragmentation. At the distinct
witness \eqref{eq:witness}, by contrast,
Proposition~\ref{prop:rejection} excludes every
substantive positive-depth union.
